\section*{Chapter \ref{ch:b3:colloids} --- Interfaces, Surfactants and Colloids}

\begin{solution}{exo:b3:colloids:1}
$r = \qty{0.50}{\micro m}$: $\Delta p = 2 \times 0.0720/\num{5.0e-7} = \qty{2.9e5}{Pa}$; inside, $1.0 + 2.9 = \qty{3.9}{bar}$.
\end{solution}

\begin{solution}{exo:b3:colloids:2}
$\ln(p/p^*) = 2 \times 0.0720 \times \num{1.807e-5}/(\num{1.0e-8} \times 8.314 \times 298.15) = 0.105$; $p/p^* = 1.11$.
\end{solution}

\begin{solution}{exo:b3:colloids:3}
$\cos\theta = (40 - 20)/72 = 0.278$, $\theta = \ang{74}$: below $90^\circ$, the liquid \omterm{def:b3:colloids:wetting}{wets} the solid
partially.
\end{solution}

\begin{solution}{exo:b3:colloids:4}
NaCl: $I = \qty{10}{mol.m^{-3}}$, $\kappa^{-1} = 0.96\sqrt{100/10} = \qty{3.0}{nm}$. \ce{MgSO4}: $I = \qty{40}{mol.m^{-3}}$, \qty{1.5}{nm}.
\end{solution}

\begin{solution}{exo:b3:colloids:5}
$\Gamma = \num{12.6e-3}/(2 \times 8.314 \times 298.15) = \qty{2.54e-6}{mol.m^{-2}}$; area $1/(\Gamma N_A) =
\qty{0.65}{nm^2}$ per molecule.
\end{solution}

\begin{solution}{exo:b3:colloids:6}
Below the CMC the slope is $-12.9$ \unit{mN/m} per unit of $\log c$ (from $-5.0$ to $-4.0$); the
plateau is at 33.4. The line reaches it at $\log c = -4.0 + (39.1 - 33.4)/12.9 = -3.56$: CMC
$\approx \qty{2.8e-4}{mol/L}$.
\end{solution}

\begin{solution}{exo:b3:colloids:7}
SDS: $P = 0.350/(0.60 \times 1.67) = 0.35$, at the border between spheres and short
cylinders: spherical \omterm{def:b3:colloids:micelle}{micelles} near the CMC. Lipid: $P = 0.700/(0.65 \times 1.67) = 0.64$:
bilayers, hence vesicles.
\end{solution}

\begin{solution}{exo:b3:colloids:8}
\ce{CaCl2}: $60/64 = \qty{0.94}{mM}$ of \ce{Ca^{2+}}; \ce{AlCl3}: $60/729 = \qty{0.082}{mM}$ of \ce{Al^{3+}}. \ce{Na3PO4}
coagulates through its \ce{Na+} (at about \qty{20}{mM} of salt); the phosphate, a
co-ion, may even adsorb and raise the negative charge.
\end{solution}

\begin{solution}{exo:b3:colloids:9}
Hydrophilic part $6 \times 44.0 + 17.0 = \qty{281.0}{g/mol}$ of \qty{450.0}{g/mol} (book atomic weights):
HLB $= 20 \times 281.0/450.0 = 12.5$, an oil-in-water emulsifier.
\end{solution}

\begin{solution}{exo:b3:colloids:10}
The oil of the small droplets is more soluble in the continuous phase (by the
factor $\exp(2\gamma V_{\mathrm m}/rRT)$, ten times larger in the exponent for \qty{50}{nm} than for \qty{500}{nm}): it
diffuses to the large droplets, which grow while the small ones vanish. This
\omterm{def:b3:colloids:ripening}{Ostwald ripening} coarsens an \omterm{def:b3:colloids:colloid}{emulsion} even if no two droplets ever merge.
\end{solution}

\begin{solution}{exo:b3:colloids:11}
The meniscus is a spherical cap of radius $r/\cos\theta$: the liquid just below it is at a
pressure lower by $2\gamma\cos\theta/r$, and the column rises until $\rho gh$ compensates. For water,
$h = 2 \times 0.072/(997 \times 9.81 \times \num{1.0e-4}) = \qty{0.15}{m}$.
\end{solution}

\begin{solution}{exo:b3:colloids:12}
At \qty{100}{mM} the \omterm{def:b3:electrode-kinetics:double-layer}{Debye length} is under \qty{1}{nm} and the repulsion dies out before the
attraction: no barrier is left (beyond the \qty{50}{mM} curve). A free polymer does not
act at the surface; a grafted polymer layer repels at contact whatever the
salt, by the entropy lost when the chains overlap.
\end{solution}

\begin{solution}{pb:b3:colloids:1}
\textbf{1.} Substitutions in the crystal lattice (\ce{Al^{3+}} for \ce{Si^{4+}}, \ce{Mg^{2+}} for \ce{Al^{3+}}) leave
a permanent negative charge, compensated by exchangeable cations.
\textbf{2.} $I = \frac12(1.0 + 1.0) = \qty{1.0}{mM}$.
\textbf{3.} $\kappa^{-1} = \qty{9.6}{nm}$.
\textbf{4.} $a\kappa = 10$: the double layer is thin compared with the particle, the
condition of the Derjaguin approximation.
\textbf{5.} $v = 2 \times 1650 \times 9.81 \times (\num{1.0e-7})^2/(9 \times \num{0.89e-3}) = \qty{4.0e-8}{m/s}$, about \qty{3.5}{mm} per day.
\textbf{6.} Their \omterm{def:b3:electrode-kinetics:double-layer}{diffuse layers} repel before van der Waals attraction takes over.
\textbf{7.} $V = 2\pi\varepsilon_r\varepsilon_0a\psi^2\ln(1 + \eu^{-\kappa h}) - Aa/12h$.
\textbf{8.} About $39\,kT$.
\textbf{9.} A fraction of order $\eu^{-39} \approx 10^{-17}$ of the collisions: practically never.
\textbf{10.} It vanishes: every collision leads to contact.
\textbf{11.} A shorter \omterm{def:b3:electrode-kinetics:double-layer}{Debye length} confines the repulsion to small gaps, where the
attraction, growing as $1/h$, dominates.
\textbf{12.} The CCC varies as $z^{-6}$ of the counterion charge.
\textbf{13.} $50/64 = \qty{0.78}{mM}$.
\textbf{14.} $50/729 = \qty{0.069}{mM}$.
\textbf{15.} The cations are the counterions of the negative particles: they
accumulate in the double layer and screen it.
\textbf{16.} Sulfate is a co-ion, repelled from the particles.
\textbf{17.} \qty{0.069}{mol} of \ce{Al^{3+}} per cubic metre.
\textbf{18.} \qty{0.034}{mol} of $\ce{Al2(SO4)3}$.
\textbf{19.} $0.0343 \times 342.3 = \qty{11.7}{g}$.
\textbf{20.} \qty{11.7}{kg} per hour.
\textbf{21.} $\ce{Al^3+ + 3 HCO3- -> Al(OH)3 + 3 CO2}$.
\textbf{22.} $3 \times 0.069 = \qty{0.21}{mM}$ of the \qty{1.0}{mM}: a fifth. The remaining
hydrogencarbonate and the dissolved \ce{CO2} buffer the water: the pH falls only
modestly.
\textbf{23.} The highly charged aluminium species adsorb and reverse the charge of the
particles, which repel each other again.
\textbf{24.} \textbf{About \qty{12}{g} of aluminium sulfate per cubic metre} (\qty{11.7}{g}), before the extra
needed for hydrolysis in practice.
\end{solution}
